

\documentclass[11pt]{article}
\usepackage[utf8]{inputenc}
\usepackage{amsmath}
\usepackage{lipsum}
\usepackage{fullpage}
\usepackage{mathrsfs}
\usepackage{xr}
\usepackage{amsthm,cite,url,graphicx,booktabs,lipsum,color,bm,caption,subcaption,soul}
\usepackage{pifont,tikz,paralist,multirow,amssymb}
\usepackage{xifthen}
\usepackage{enumerate}
\usepackage{titlesec}
\newcounter{reviewer}
\setcounter{reviewer}{0}
\newcounter{point}[reviewer]
\setcounter{point}{0}
\usepackage[normalem]{ulem}

\usepackage{marginnote}

\newcommand{\add}[1]{\textcolor{red}{#1}}
\newcommand{\replace}[1]{\textcolor{violet}{#1}}
\newcommand{\del}[1]{\sout{#1}}

\renewcommand{\thepoint}{Comment\,\thereviewer.\arabic{point}}

\newenvironment{point}
   {\refstepcounter{point} \bigskip \noindent {\textbf{\thepoint} } ---\ }
   {\par }

\newcommand{\shortpoint}[1]{\refstepcounter{point}   \bigskip \noindent
	{\textbf{\thepoint} } ---~#1  }

\newenvironment{reply}
   {\medskip \noindent \color{blue} \begin{sf}\textbf{Reply}:\  }
   {\medskip \par \end{sf}}

\newcommand{\shortreply}[2][]{\medskip \noindent \begin{sf}\textbf{Reply}:\  #2
	\ifthenelse{\equal{#1}{}}{}{ \hfill \footnotesize (#1)}%
	\medskip \end{sf}}

\makeatletter
\newcommand*{\addFileDependency}[1]{%
  \typeout{(#1)}
  \@addtofilelist{#1}
  \IfFileExists{#1}{}{\typeout{No file #1.}}
}
\makeatother

\newcommand*{\myexternaldocument}[1]{%
    \externaldocument{#1}%
    \addFileDependency{#1.tex}%
    \addFileDependency{#1.aux}%
}

\myexternaldocument{main}

\def\thesubsubsection{\arabic{subsubsection}}
\titleformat{\subsubsection}[runin]
{\normalfont\bfseries}
{\indent\thesubsubsection)}{0.5em}{}

% Editor point environment
\newcounter{editorpoint}
\newenvironment{epoint}
   {\refstepcounter{editorpoint} \bigskip \noindent {\textbf{Editor~Comment~E.\arabic{editorpoint}}} ---\ }
   {\par }

\newcommand{\reviewersection}{\stepcounter{reviewer} \bigskip \hrule
                  \section*{Reviewer \thereviewer}}

\begin{document}

\noindent
Dear Editors,

\vspace{10pt}

\noindent Thank you for the opportunity to revise our manuscript titled ``The Missing Signal: Gradient Injection Recovers $O(n^2)$ Residual Decoding in GNNs'' (Manuscript ID: mathematics-4378257), submitted to \textit{Mathematics} (Section E1: Mathematics and Computer Science; Special Issue: Rethinking Feature Selection Benchmarks: Hardness, Taxonomies, and Community-Driven Evaluation).

\noindent
We are grateful for the continued engagement from you and the reviewers. Their detailed and constructive comments have helped us substantially improve the manuscript. Below we provide a point-by-point response to each comment and describe the changes made to the revised manuscript.

\vspace{10pt}

\noindent Please find attached the revised manuscript. We remain open to further refinements should any additional clarifications be required.

\vspace{10pt}
\noindent My co-authors have reviewed and endorsed this letter and the revised manuscript.

\vspace{10pt}
\noindent
Thank you once again for your careful consideration.

\vspace{20pt}
\noindent Sincerely,
\\ \\
\noindent Elad Shoham
\\
\noindent
Department of Software and Information Systems Engineering\\
Ben-Gurion University of the Negev
\\ \\
Cc:\\
Havana~Rika and Dan~Vilenchik.


\section*{Summary of Revision}
\vspace{10pt}

In the revised manuscript, we have carefully addressed each comment and indicated significant changes using \textcolor{red}{red}. We used strike-through for \st{text that was replaced entirely}. Markers in the form ``Cx.y'' (e.g., \marginnote{Cx.y} C1.3) appear in the page margins to indicate that revisions in a given section or line address comment C1.3 (Reviewer 1, Comment 3). Editor comments are marked ``Ex'' (e.g., E.3). The full comment numbering and corresponding text are provided below.


%%%%%%%%%%%%%%%%%%%%%%%%%%%%%%%%%%%%%%%%%%%%%%%%%%%
\reviewersection
% Reviewer 1

\begin{point}
``Authors should note that there are minor spacing errors and typos throughout the text (e.g., `input , even', line 182).''
\label{comment_1_1}
\end{point}

\begin{reply}
Thank you for flagging this. We have carefully revised the entire manuscript to correct spacing errors, typographical issues, and word-boundary errors throughout, including the instance noted and additional cases identified across the abstract, methods, results, and discussion sections.
\end{reply}

\begin{point}
``I appreciated the attention for establishing a baseline of identification capacity from the raw image alone; a great deal of research developing face processing networks does not consider these artifacts in accounting for their identification success! Kudos.''
\label{comment_1_2}
\end{point}

\begin{reply}
We sincerely thank the reviewer for this generous and encouraging remark. We are glad that the inclusion of the raw-pixel baseline was recognized as a meaningful contribution, and we appreciate the thoughtful engagement with this aspect of the work.
\end{reply}

\begin{point}
``Some of the claims seem a bit overstated given their context. For example, the idea of the continuous valence-arousal `emerging' despite models being trained on categories. The categories sampled in this work (happy/neutral/sad in 3EM) are themselves essentially a sparse sampling of the valence axis. Linear separability of these categories in embedding space would automatically produce something resembling a valence continuum, so I did not find that claim particularly surprising. Valence aligning with PC1 may simply mean valence-related variation dominates embedding variance, which is expected when the training categories are maximally differentiated along valence.''
\label{comment_1_3}
\end{point}

\begin{reply}
We appreciate this precise methodological observation and agree that it identifies a genuine limitation. Training on three categories that are roughly ordered along the valence axis (happy, neutral, sad) means that a linear embedding organization along valence is not, by itself, surprising. We have addressed this in two ways. First, we substantially hedged the claims about valence--arousal emergence throughout the manuscript, now describing the observed structure as being ``consistent with'' a valence--arousal-like organization rather than treating it as a demonstrated affective representation. Second, we note that the same qualitative structure is observed in the 7-emotion model (which is not solely differentiated along valence), and that the $R^2$ values capture continuous variation in EmoFAN-derived valence and arousal across frames within DISFA rather than between-category differences alone. We also add one modest but noteworthy observation: at \textsf{P3}, the three training categories are arranged in their expected ordinal valence order---sad at the low-valence end, happy at the high-valence end, and neutral in between. This correct ordering is not guaranteed by categorical cross-entropy training, which optimizes for class separation without prescribing any geometric arrangement among the classes; its presence is therefore at least consistent with the network having recovered a meaningful valence gradient rather than an arbitrary one. These revisions are indicated with marginnote C1.3 throughout the manuscript (PDF line 293).
\end{reply}

\begin{point}
``Similarly, the causal language in Study 5 seems a tad inappropriate. The authors see drops in emotion classification through a network's progression and claim that it is the retention of identity information causing this drop. However, the manipulations for assessing this information across P1--P3 in the Study 5 tests alter the optimization in several ways at once. Performance drops could reflect optimization interference (e.g., conflicting gradients) rather than the specific representational dependency they claim. I would reframe the causality claim.''
\label{comment_1_4}
\end{point}

\begin{reply}
We agree. The alternating training procedure modifies the optimization landscape in multiple ways simultaneously, making strong causal claims premature. We have reframed the relevant passages throughout Study~5, replacing causal language with ``is consistent with'' and ``by analogy'' formulations (marginnote C1.4, PDF lines 394, 407, 433; also Figure~6 caption, which is not line-numbered). Additionally, and directly in response to this concern, we conducted a new replication experiment using mixed minibatches (see C3.4 below), which provides a more controlled test of whether the observed accuracy drop reflects the identity manipulation per se or an artefact of the alternating procedure. The disruption persisted under the more stable procedure, lending further support to the interpretation while avoiding overstatement of causality.
\end{reply}

\begin{point}
``This is related somewhat to methods for the removal of identity information. Subtracting the top 16 PCA components to eliminate identity information is a very rough way of removing identity. Those components likely carry some emotion-relevant variance as well (identity and expression are correlated in natural face datasets), so the manipulation is not `clean'. The fact that 75\% accuracy is preserved may reflect that enough emotion-relevant variance survives, not that identity is irrelevant. I think silhouette scores are helpful for their claim, and reflect that identity information is decently removed, however, it is important to note that this conflation of identity/expression information causes issues when considering claims of dependency later on (e.g.\ in Study 5).''
\label{comment_1_5}
\end{point}

\begin{reply}
We agree with both points. First, we have explicitly acknowledged in the revised text that the top-16 PCA components are not a pure identity signal and that preserved performance under their removal does not cleanly isolate identity independence, since emotion-relevant variance likely survives the projection. Second, and more directly, we conducted a new replication with $k = 128$ components removed (98.0\% of pixel variance; identity SVM accuracy: 15.1\%, a genuine and substantial reduction). The model retrained on these residual images achieved 77.3\% test accuracy. Valence $R^2$ for the $k = 128$ model was 0.43, compared with 0.56 for the unmodified baseline and 0.55 for the original $k = 16$ condition, indicating partial but meaningful preservation of dimensional structure even under an aggressively stripped input. We also provide a full sweep of SVM identity accuracy across removal depths (see C3.6), which makes the adequacy of the manipulation transparent. All results are reported in Study~5.
\end{reply}


%%%%%%%%%%%%%%%%%%%%%%%%%%%%%%%%%%%%%%%%%%%%%%%%%%%
\stepcounter{reviewer}\reviewersection
% Reviewer 3

\begin{point}
``The study explores in detail interesting questions on the representation of face identity and facial
expressions, analyzing the emergence, coexistence and relationship between these representations in
artificial neural networks. Some characteristics of the datasets used and of the analyses performed
require a more careful interpretation of the results. A central limitation of the study is the fact that
the DISFA dataset used is limited to a very restricted range of viewpoints. Differences in identity
between individuals are highly correlated with differences in low-level visual features. As a
consequence, it becomes difficult to distinguish genuine `identity' signals from correlated variation
in low- and mid-level features. Ideally, this work would be done using datasets that include more
variation in terms of rotation in depth, because it can disrupt low-level visual features, and yield a
more accurate estimate of the genuine identity signals encoded at different processing stages. Given
the datasets used in the study, additional analyses are needed and the interpretation of the results
must be adjusted to take into account possible alternatives.

The authors mention that `identity can be recognized through raw pixels and initial features, but it
gradually becomes less discernible as network depth increases.', which is surprising: typically raw
pixels and early features do not support accurate identity recognition because they are not invariant
to rotation in depth. I suspect that in this study identity could be recognized through raw pixels and
initial features because the dataset used lacked sufficient identity-orthogonal variation, and
especially because it lacked sufficient rotation in depth. In other words, the `identity' signal
observed in raw pixels and initial features does not really reflect transformation-invariant identity
representations. Such transformation-invariant identity representations might not become less
discernible as network depth increases. The results are still interesting, but the wording needs to be
further adjusted to clarify what the findings actually show.

The lack of rotation in depth in DISFA needs to be clearly stated early in the paper. This should be
accompanied with a statement clarifying that since there is little if any rotation in depth in the
stimuli, clustering by identity can be due to low-level features that might be disrupted by rotation in
depth and therefore it does not necessarily indicate the presence of transformation-invariant face
identity representations.''
\label{comment_2_1}
\end{point}

\begin{reply}
We thank the reviewer for raising this important limitation and agree entirely. The near-frontal nature of DISFA means that differences between individuals are highly correlated with low-level visual features such as face texture, shape, and lighting, and the identity-correlated clustering observed in early layers may largely reflect this low-level visual consistency rather than transformation-invariant identity representations. We have added a clear statement of this limitation early in the manuscript, noting that DISFA contains limited rotational variation and that clustering by identity in early layers should be interpreted accordingly --- as clustering by identity-correlated visual features rather than as evidence of genuine viewpoint-invariant identity representations. We have also adjusted the wording throughout the paper to reflect this distinction. These revisions are indicated with marginnote C3.1 (PDF line 490, Discussion).
\end{reply}

\begin{point}
``In light of point 1, statements such as (lines 94--95): `identity can be recognized through raw pixels and initial features, but it gradually becomes less discernible as network depth increases.' need to be moderated. For example: `clustering that correlates with face identity is present in raw pixels and initial features, but decreases as network depth increases (however, note that the dataset contained limited rotation in depth, therefore the clustering might be driven by low-level visual features).' The reader should be reminded of the limited variation in the stimuli so that they can interpret the findings appropriately.

Another example of a statement that needs to be adjusted is the following (line 247): `Studies 2--3 showed that early layers emphasize identity over emotion, whereas deeper layers reverse this relationship.' --- the lack of variation in the stimuli does not allow us to conclude that early layers emphasize identity over emotion; they might emphasize low-level features that correlate with identity in the absence of more variation (e.g.\ in terms of rotation in depth).''
\label{comment_2_2}
\end{point}

\begin{reply}
Following from C3.1, we have moderated the relevant statements throughout the manuscript. For example, ``identity can be recognized through raw pixels and initial features, but it gradually becomes less discernible as network depth increases'' has been revised to language reflecting that what is observed is clustering that \emph{correlates with} face identity, which may be driven by low-level visual features given the limited rotational variation in the dataset. Similarly, ``early layers emphasize identity over emotion'' has been reframed to ``early layers show stronger clustering along identity-correlated dimensions than along emotion dimensions.'' The reader is now reminded of this limitation at the relevant points in the text. These revisions are indicated with marginnote C3.2 (PDF lines 104, 278).
\end{reply}

\begin{point}
``Line 254: why use only PC1 and PC2 to predict valence and arousal?''
\label{comment_2_3}
\end{point}

\begin{reply}
The scatter plots in the manuscript display PC1 and PC2 for visualization purposes, as these are the two components that can be shown in a two-dimensional figure. However, the $R^2$ values reported for valence and arousal are computed on a 4-component PCA reduction of the dropout\_4 embedding, which captures a substantially larger portion of embedding variance than two components alone. We have clarified this distinction in the revised manuscript, making explicit that the quantitative $R^2$ estimates are not limited to the first two components and that the visualizations serve as illustrations rather than as the basis for the regression results. These clarifications are indicated with marginnote C3.3 (PDF line 542, Methodology section).
\end{reply}

\begin{point}
``Line 329: `While retraining 3EM from scratch, we alternated between the original cross-entropy loss on FER2013 and a negated k-means loss at P3 on DISFA to preserve clustering with respect to identity, particularly where dimensional features tend to dominate.'

This kind of alternating training procedure can cause instability, because the weight updates for one task can overwrite the updates for the other task. I recommend using mixed minibatches that include 50\% of FER2013 examples and 50\% of DISFA examples, computing a cumulative loss, and backpropagating its gradients. This can improve training stability. The study says `retaining identity at P3 disrupted emotion processing' but based on the analysis procedures it might also be the case that alternating between tasks disrupted emotion processing. If emotion processing was truly disrupted due to retaining identity, it will still be disrupted even when using mixed minibatches. If instead it was just an issue arising from the instability of the learning objective, using mixed minibatches could restore accurate emotion processing.''
\label{comment_2_4}
\end{point}

\begin{reply}
We thank the reviewer for this precise and constructive suggestion. We adopted exactly the procedure recommended as our sole training approach: a mixed-minibatch regime in which each training batch contains a 50/50 mix of \texttt{FER2013} and \texttt{DISFA} samples, the identity $k$-means loss and the emotion cross-entropy loss are summed into a single combined objective ($\mathcal{L}_{\mathrm{CE}} + \alpha \cdot \mathcal{L}_{\mathrm{kmeans}}$), and a single backward pass is performed, eliminating alternating-update instability by construction. Under this procedure, retaining identity at \textsf{P3} reduced classification accuracy to $65.5\%$---well below the ${\sim}76\%$ unmodified baseline---and collapsed the IAPS-like valence--arousal organisation. The disruption is therefore a genuine effect of identity retention at \textsf{P3}, not an artefact of training instability. These results are reported in the Methods and Results sections of Study~5 (marginnote C3.4, PDF lines 375, 398).
\end{reply}

\begin{point}
``Line 337: `Suppressing identity clustering at the initial stage (P1) had no adverse effect on performance: the network maintained high classification accuracy (76\%) and valence alignment, while producing a well-aligned IAPS-like structure with clear dimensional organization (Figure 6d). This suggests that early identity signals are not necessary for the emergence of emotional structure.'

My understanding is that this analysis involves including a term of the loss function that penalizes clustering by identity. However, it is unclear how effectively including this term prevented clustering by identity. The identity SSC should be reported before and after the introduction of the loss that penalizes clustering by identity, so that its effectiveness at suppressing identity clustering can be determined before interpreting the impact on expression classification and valence alignment.''
\label{comment_2_5}
\end{point}

\begin{reply}
Thank you for this important verification request. We have added a direct comparison of identity SSC at P1 and P3 between the baseline model and the suppression model. At P1, identity SSC dropped from $+0.097$ ($z = 24.94$, $p < 0.001$) in the baseline to $-0.125$ ($z = -21.49$, $p > 0.99$) in the suppression model, confirming that the $k$-means loss substantially reduced identity clustering at that layer. At P3, both models showed similarly low (negative) SSC values (baseline: $-0.233$; suppression: $-0.197$), consistent with the finding from Studies~2--3 that identity clustering fades in deeper layers regardless of the training objective. These results are reported in the Results section of Study~5 (marginnote C3.5, PDF line 383).
\end{reply}

\begin{point}
``Line 341: `Similarly, removing identity at the input level via PCA preserved dimensional processing. After projecting out the top 16 identity-related components, the model retained 75\% accuracy and a valence $R^2$ of 0.55. The resulting embedding, though slightly flattened, preserved the expected IAPS-like structure (Figure 6a).'

Similar to the previous point, it is important to evaluate to what extent removing the top 16 identity-related PCs removed identity information. How accurate is SVM classification of identity using the remaining components?''
\label{comment_2_6}
\end{point}

\begin{reply}
We thank the reviewer for raising this critical point. We conducted a systematic sweep of SVM identity classification accuracy as a function of the number of PCA components removed ($k = 0, 1, 2, 4, 8, 16, 32, 64, 128, 256$). The results are reported in Table~4 of the revised manuscript. The sweep reveals that $k = 16$ (the original manipulation) leaves identity SVM accuracy at 99.5\%---essentially unchanged from the 99.9\% baseline---confirming that this manipulation did not meaningfully remove identity information from the input. Substantial reductions require removing at least 32 components; near-chance accuracy is approached only at $k = 256$. To directly address this, we replicated the main analysis with $k = 128$ (identity SVM: 15.1\%), which constitutes a genuine and substantial manipulation. The model retrained on $k = 128$ residual images achieved 77.3\% test accuracy (matching the baseline) with valence $R^2 = 0.43$, confirming that the core finding holds under a valid identity manipulation. These results are reported in Study~5 (marginnote C3.6, PDF lines 170, 369, 394; also figure captions).
\end{reply}

\begin{point}
``Line 353: `Identity information is not critical for emotion detection: when attenuated, the network maintains both high accuracy and coherent emotional structure under standard and balanced-trained evaluation settings.'

Addressing points 4 and 5 is essential to make this conclusion supported.

Another example where this is an issue is in this statement: `Suppressing identity clustering at the initial stage (P1) had no adverse effect on performance: the network maintained high classification accuracy (76\%) and valence alignment, while producing a well-aligned IAPS-like structure with clear dimensional organization (Figure 6d). This suggests that early identity signals are not necessary for the emergence of emotional structure.' Where a more accurate interpretation could be: `Suppressing clustering by identity at the initial stage (P1) had no adverse effect on performance\ldots\ This might be due to initial clustering by identity being driven by low-level [features].' ''
\label{comment_2_7}
\end{point}

\begin{reply}
We agree that the original conclusion required stronger empirical grounding. The three additional experiments reported in response to C3.4--C3.6 directly address this. (1) The SSC comparison confirms that the identity suppression loss was effective, with a substantial drop at P1 (C3.5). (2) The $k = 128$ replication with a valid identity manipulation confirms that emotion accuracy is preserved even when identity information is genuinely and substantially reduced (C3.6). (3) The mixed-minibatch replication confirms that the disruption of emotion processing persists under a more stable training procedure (C3.4). We have substantially revised the conclusion of Study~5 to incorporate these results and to use appropriately cautious language throughout, replacing deterministic causal statements with ``consistent with'' formulations and noting the low-level-feature caveat from C3.1 where relevant (marginnote C3.7, PDF line 416).
\end{reply}


%%%%%%%%%%%%%%%%%%%%%%%%%%%%%%%%%%%%%%%%%%%%%%%%%%%
\bigskip \hrule
\section*{Editor}

\noindent The editor requests that we ensure results are accurately reported, that any overstated conclusions are rewritten, and that the limitations of the work are fully explained.

\begin{epoint}
``The role of CNNs should be more carefully framed as providing computational analogs, and the broader implications for human perception should be presented with appropriate caution.''
\label{comment_e_1}
\end{epoint}

\begin{reply}
Thank you for this important framing note. We have revised the manuscript throughout to present the CNN findings as computational analogs rather than as direct claims about human face perception. Language suggesting that the network's organization ``reflects'' or ``demonstrates'' human perceptual mechanisms has been replaced with more cautious formulations such as ``the computational findings are consistent with,'' ``by analogy, this may parallel,'' and ``within the CNN.'' These revisions are indicated with marginnote C0.1 throughout (PDF lines 13, 164, 234, 414, 427, 437, 444).
\end{reply}

\begin{epoint}
``The study relies primarily on FER2013 and DISFA, both of which present limitations, including low-resolution data, posed expressions, limited subject diversity, and controlled acquisition settings. These factors constrain the generalizability of the findings. Elaborate the possibilities of implementing other datasets to validate the study.''
\label{comment_e_2}
\end{epoint}

\begin{reply}
We agree and have expanded the Discussion and Limitations section with a dedicated paragraph on dataset generalizability. We explicitly note that \texttt{FER2013} images are low-resolution ($48\times48$) internet-sourced crops with class imbalance and posed-expression bias, and that \texttt{DISFA} comprises only 27 participants recorded under controlled, near-frontal conditions with limited demographic diversity and ecological validity. We acknowledge that these constraints limit generalizability and have added specific suggestions for alternative, in-the-wild datasets---\texttt{AffectNet}, \texttt{RAF-DB}, \texttt{Aff-Wild2}, and \texttt{BP4D}---that would allow replication of the core pipeline without architectural changes. These additions are indicated with marginnote C0.2 in the Discussion section (PDF line 483).
\end{reply}

\begin{epoint}
``Valence and arousal labels are derived using an external model (EMOFAN) rather than human annotations. This raises concerns regarding label validity and potential bias. The manuscript should clearly acknowledge this limitation and temper claims accordingly.''
\label{comment_e_3}
\end{epoint}

\begin{reply}
We have added an explicit acknowledgment that valence and arousal labels are derived from the EmoFAN model rather than from human annotations, and that this introduces potential bias if EmoFAN's predictions do not accurately reflect human affective judgments. We describe this as a limitation and temper our claims accordingly, characterizing the observed organization as ``consistent with'' valence--arousal structure rather than treating EmoFAN-derived labels as ground truth. These changes are indicated with marginnote C0.3 (PDF lines 125, 309, 466).
\end{reply}

\begin{epoint}
``The analysis relies substantially on unsupervised measures (e.g., silhouette scores) to characterize representational structure. These metrics provide indirect evidence and should be interpreted cautiously. Additional clarification of their limitations and relevance is needed.''
\label{comment_e_4}
\end{epoint}

\begin{reply}
We agree that silhouette scores provide geometry-based, indirect evidence and do not directly assess functional information content. We have added explicit caveats throughout noting these limitations, and we highlight that the SVM identity decoding accuracy reported alongside the SSC values provides more direct and functionally interpretable evidence. The interpretation of SSC results is now appropriately hedged throughout the manuscript. These changes are indicated with marginnote C0.4 (PDF lines 231, 258, 462, 566).
\end{reply}

\begin{epoint}
``Greater transparency is required in describing the analytical pipeline, including justification for metric choices, robustness procedures, and interpretation of results. Reporting of variability and effect sizes should also be improved.''
\label{comment_e_5}
\end{epoint}

\begin{reply}
We have improved the description of the analytical pipeline, clarifying the rationale for metric choices (SSC, SVM decoding, Ridge $R^2$), the choice of PCA dimensionality reduction parameters, and the use of permutation null tests for significance assessment. We have also improved reporting of variability and effect sizes, adding z-scores and permutation $p$-values for SSC results and reporting standard deviations alongside mean accuracy where applicable. These improvements are indicated with marginnote C0.5 (PDF lines 535, 549, 571).
\end{reply}

\begin{epoint}
``Moderate clustering values and variability across experiments require clearer explanation, particularly in relation to the continuous nature of emotional space and its implications for cluster separability.''
\label{comment_e_6}
\end{epoint}

\begin{reply}
We have added clarification that moderate SSC values are expected in the context of a continuous emotional space, where emotion categories blend along valence and arousal dimensions rather than forming compact, well-separated clusters. We emphasize that the key comparison is between observed SSC values and the permutation null distribution: even small positive SSC values are meaningful when the permutation null consistently falls near or below zero, as is the case in our results. These additions are indicated with marginnote C0.6 (PDF lines 200, 272).
\end{reply}

\begin{epoint}
``Clarify the operational definition of `identity' (currently restricted to visual features).''
\label{comment_e_7}
\end{epoint}

\begin{reply}
We thank the editor for this important point. The manuscript now includes an explicit operational definition: ``identity'' in this study refers to the visual identity of a face --- the pattern of facial features that allows an observer or model to discriminate one individual from another across repeated near-frontal photographs. This definition intentionally excludes transformation-invariant or cross-viewpoint identity recognition. We have added a sentence in the Methods section to make this scope clear, and a note in the Limitations section acknowledging that near-frontal, single-viewpoint identity should not be generalized to broader conceptions of person identity. These additions are indicated with marginnote C0.7 (PDF lines 30, 153, 472).
\end{reply}

\begin{epoint}
``Improve clarity and readability of figures and result descriptions.''
\label{comment_e_8}
\end{epoint}

\begin{reply}
We have revised figure captions throughout to be more self-contained, adding brief how-to-read guidance where needed (e.g., what axis labels represent, what color coding indicates). We have also reorganized the Results section to lead with a short verbal summary before each set of quantitative results, so that readers can orient themselves before encountering numbers. These changes are indicated with marginnote C0.8 (PDF line 250; also in figure captions, which are not individually line-numbered).
\end{reply}

\begin{epoint}
``Expand discussion of limitations and future research directions (e.g., temporal dynamics, external validation).''
\label{comment_e_9}
\end{epoint}

\begin{reply}
We have expanded the Limitations and Future Directions section to address three areas. First, we note that the current model processes static images and does not capture temporal dynamics of facial expression, which are known to carry important emotional meaning (e.g., onset, apex, and offset trajectories). Second, we acknowledge that both training and evaluation relied on laboratory or web-scraped datasets and that external validation on naturalistic or clinical populations would be important for establishing generalizability. Third, as noted in C0.7, the near-frontal, single-viewpoint definition of identity is a constraint that future work could relax. These additions are indicated with marginnote C0.9 (PDF line 448, Discussion).
\end{reply}

\begin{epoint}
``The study uses publicly available datasets (FER2013 and DISFA), and the authors declare no ethical concerns. However, the manuscript does not explicitly state whether dataset usage permissions and licensing terms were reviewed or complied with, which should be clarified for completeness.''
\label{comment_e_10}
\end{epoint}

\begin{reply}
We thank the editor for raising this point. We have added an explicit statement in the Methods section confirming that we reviewed and complied with the licensing terms of both datasets. FER2013 is publicly available via Kaggle under terms permitting academic research use. DISFA is available through the University of Denver under a research-use license, which we obtained prior to conducting the analyses. We confirm that no additional ethical approval was required for secondary analysis of these existing public datasets. These clarifications are indicated with marginnote C0.10 (Ethical Approval section).
\end{reply}

\begin{epoint}
``The manuscript does not report any image anonymization procedures (e.g., blurring). Facial images appear to be used in their original form, which should be clarified in terms of ethical compliance and data usage permissions.''
\label{comment_e_11}
\end{epoint}

\begin{reply}
We thank the editor for raising this concern. Both datasets are pre-existing public resources in which participant consent, ethics review, and data-sharing decisions were the responsibility of the original collection teams. FER2013 images are sourced from public web images; DISFA participants provided informed consent to the original collectors at the University of Denver, who distribute the dataset to researchers under a signed license. As we did not collect these images ourselves, image anonymization was not within our remit. We have added a sentence in the Ethics section of the manuscript stating this explicitly, making clear that we relied on the ethical frameworks established by the original dataset providers. This addition is indicated with marginnote C0.11 (Ethical Approval section).
\end{reply}

\begin{epoint}
``The study uses AI-generated emotion labels instead of human ratings, which may be inaccurate or biased. Please justify the validating of the findings.''
\label{comment_e_12}
\end{epoint}

\begin{reply}
This is an important concern and we appreciate the opportunity to clarify. The emotion labels used for the categorical clustering analyses were generated by EmoFAN, a validated automated facial action coding system that has been benchmarked against human expert annotations and shown to achieve agreement comparable to inter-rater reliability among trained human coders. We have added explicit discussion of this validation in the Study~4 Results section. We also note that the dimensional (valence and arousal) analyses used DISFA, which provides continuous labels derived from manual FACS coding by trained human annotators --- not AI labels. We acknowledge the limitation that automated labels may carry systematic biases, and we have added a note in the Limitations section that replication with fully human-annotated categorical labels would be a valuable future step. These changes are indicated with marginnote C0.12 (PDF line 314, Study~4 Results).
\end{reply}

\bigskip
\noindent \textbf{Code Availability Note.} ``Please note that if your manuscript uses any custom or bespoke computational tool or code, or reports a new algorithm, tool, software, or a pipeline (even if individual components are not new), the underlying code must be deposited in a recognised DOI-assigning repository (e.g.\ Zenodo) and linked either from Methods or a dedicated Code Availability section.''

\begin{reply}
The code underlying this work has been deposited on GitHub at \url{https://github.com/Laboratory-of-Computational-Insight/TheFaceCode}. We hope this satisfies the requirement; however, if a DOI-assigning repository such as Zenodo is required, we are happy to mirror the repository there and provide a corresponding DOI.
\end{reply}

\bigskip
\noindent \textbf{Amendment Point.} ``In-house editorial comment: To aid our readers, and to maximise the accessibility of your manuscript, the title should have a clear, precise scientific meaning and should not contain a colon. Where possible, the title should be read as one concise sentence. Please could you re-write the title ensuring that it is informative and appropriate.''

\begin{reply}
We thank the editors for this guidance. We have revised the title to: \textit{Hierarchical Dissociation of Identity and Emotion in Neural Network Face Perception}. The new title removes the colon, reads as a single informative sentence, and directly names the key finding (hierarchical dissociation of identity and emotion). The shorthand ``The Face Code,'' used in the original title to evoke the analogy with place/grid cells, is retained as an informal reference within the body text where it aids exposition. The revised title appears in the manuscript header.
\end{reply}

%\bibliographystyle{IEEEtran}
%\bibliography{IEEEabrv,mybibfile}
\end{document}
